\documentclass[11pt]{article}
\usepackage[a4paper,margin=2.4cm]{geometry}
\usepackage{amsmath,amssymb}
\usepackage{enumitem}
\usepackage{hyperref}

\newcommand{\lam}[1]{\lambda #1.\,}
\newcommand{\Lam}[1]{\Lambda #1.\,}
\newcommand{\List}{\mathsf{List}}
\newcommand{\nat}{\mathsf{nat}}

\title{Lecture 3 In-Class Worksheet\\Polymorphic Lambda Calculus}
\author{}
\date{}

\begin{document}
\maketitle

\section*{Learning objectives}
By the end of this three-hour unit, students should be able to:
\begin{enumerate}[nosep]
  \item explain the difference between term abstraction and type abstraction;
  \item typecheck simple System F terms;
  \item instantiate polymorphic functions using type application;
  \item read Church encodings of sums, products, naturals, and lists as eliminators;
  \item state the basic intuition of parametricity without proving the full theorem.
\end{enumerate}

\section*{Suggested pacing}
\begin{center}
\begin{tabular}{r p{0.68\linewidth}}
0--20 min & From STLC to explicit polymorphism.\\
20--50 min & System F syntax and typing rules.\\
50--75 min & Exercise 1: type derivations.\\
75--85 min & Break.\\
85--125 min & Programming with polymorphic encodings.\\
125--150 min & Exercise 2: products and sums.\\
150--175 min & Parametricity as a reasoning principle.\\
175--180 min & Exit check.\\
\end{tabular}
\end{center}

\section*{Exercise 1: polymorphic derivations}
Derive or explain the following judgments.
\begin{enumerate}
  \item $\vdash \Lam{X}\lam{x:X}x : \forall X. X \to X$
  \item $\vdash \Lam{X}\Lam{Y}\lam{x:X}\lam{y:Y}x : \forall X.\forall Y. X \to Y \to X$
  \item Why is $\Lam{X}\lam{x:X}x\;X$ ill-formed as written?
\end{enumerate}

\paragraph{Solution checkpoints.}
\begin{enumerate}[nosep]
  \item Introduce type variable $X$, derive $x:X \vdash x:X$, abstract over $x$, then over $X$.
  \item Same pattern with two type variables and two term variables.
  \item Type application applies a term to a type only when the term has a universal type; here $x : X$, not $\forall Y.A$.
\end{enumerate}

\section*{Exercise 2: Church encodings as eliminators}
Use the encodings
\[
  A \times B \equiv \forall X. (A \to B \to X) \to X,
  \qquad
  A + B \equiv \forall X. (A \to X) \to (B \to X) \to X.
\]
Define:
\begin{enumerate}
  \item $\mathsf{proj}_1 : A \times B \to A$;
  \item $\mathsf{proj}_2 : A \times B \to B$;
  \item $\mathsf{either} : \forall X. (A \to X) \to (B \to X) \to A + B \to X$.
\end{enumerate}

\paragraph{Solution checkpoints.}
\begin{enumerate}[nosep]
  \item $\mathsf{proj}_1\;p \equiv p\;A\;(\lam{x}\lam{y}x)$.
  \item $\mathsf{proj}_2\;p \equiv p\;B\;(\lam{x}\lam{y}y)$.
  \item $\mathsf{either}\;X\;f\;g\;s \equiv s\;X\;f\;g$.
\end{enumerate}

\section*{Exercise 3: parametricity intuition}
For each type, list the closed terms students should expect before seeing any implementation.
\begin{enumerate}
  \item $\forall X. X$
  \item $\forall X. X \to X$
  \item $\forall X. X \to \nat$
\end{enumerate}

\paragraph{Solution checkpoints.}
\begin{enumerate}[nosep]
  \item No closed total term in pure System F.
  \item The identity function is the canonical inhabitant.
  \item Such a term cannot inspect an arbitrary $X$; it must return a fixed natural number.
\end{enumerate}

\section*{Exit check}
Students should be able to answer these without notes:
\begin{enumerate}[nosep]
  \item What is the difference between $\lambda x.t$ and $\Lambda X.t$?
  \item What does universal elimination do?
  \item Why does $\forall X. X \to X$ strongly suggest the identity function?
\end{enumerate}

\section*{Instructor notes}
For a no-homework version, choose one main application: either Church encodings or parametricity. If both are included, keep parametricity at the level of examples and free-theorem intuition, not the full logical-relations proof.

\end{document}
